\chapter[Band 27: Endliche Wörter und Klammerungsbäume]{Band 27 auf einen Blick}
\noindent\textbf{Endliche Wörter und Klammerungsbäume}\par\medskip

% Dieselbe lokale Notation wie im Fachband, für die übernommene Übersicht.
\begingroup
% Abstand zur Tabellenlinie auch unter mehrzeiligen Formeln sichern.
\setlength{\aboverulesep}{6pt}
\newcommand{\BinaryAddresses}{\mathcal W_2}
\newcommand{\AddressPrefixMap}[1]{\operatorname{vor}_{#1}}
\newcommand{\AddressPrefix}[2]{\AddressPrefixMap{#1}(#2)}
\newcommand{\AddressGraft}[2]{\Gamma(#1,#2)}
\newcommand{\GoodAddressSet}[1]{\mathsf{VollPraef}(#1)}
\newcommand{\TreePositions}[2]{\operatorname{Pos}_{#1}(#2)}
\newcommand{\AddressEdges}[1]{E_{#1}^{\mathrm{adr}}}
\newcommand{\AddressInner}[1]{I_{#1}^{\mathrm{adr}}}
\newcommand{\AddressLeft}[1]{\lambda_{#1}^{\mathrm{adr}}}
\newcommand{\AddressRight}[1]{\rho_{#1}^{\mathrm{adr}}}
\newcommand{\TreePositionGraph}[2]{\mathcal G_{#1}^{\mathrm{adr}}(#2)}
\begingroup
\small
\setlength{\LTleft}{0pt}
\setlength{\LTright}{0pt}
\setlength{\LTpre}{.5\baselineskip}
\setlength{\LTpost}{0pt}
\renewcommand{\arraystretch}{1.08}
\newcommand{\BXXVIIBlickFormel}[1]{%
  \par\smallskip\noindent\(\displaystyle #1\)\par}
\newcommand{\BXXVIIBlickQuellen}[1]{%
  \par\nobreak\smallskip{\footnotesize #1}\par}

\noindent
\emph{Lesart.}
\(\NatSeg n=\{k\in\mathbb N\mid k<n\}\) ist der \(n\)-te Anfangsabschnitt.
Links von \(\vdash\) stehen die Voraussetzungen, rechts die Folgerung;
\(\eqvdash\) bezeichnet die entsprechende Äquivalenz.
\(\mathsf{BinOp}(A,\star)\) bedeutet, dass \(\star\) eine binäre Operation
auf \(A\) ist.

\begin{longtable}{@{}>{\raggedright\arraybackslash}p{.455\linewidth}%
  @{\hspace{.012\linewidth}\vrule width .35pt\hspace{.012\linewidth}}%
  >{\raggedright\arraybackslash}p{.51\linewidth}@{}}
\toprule
\textbf{Definitionen und Konstruktionen}&\textbf{Zentrale Resultate}\\
\midrule
\endfirsthead
\toprule
\textbf{Definitionen und Konstruktionen}&
  \textbf{Zentrale Resultate (Fortsetzung)}\\
\midrule
\endhead
\midrule
\multicolumn{2}{@{}r@{}}{%
  \footnotesize\itshape Fortsetzung auf der nächsten Seite}\\
\endfoot
\bottomrule
\endlastfoot

\multicolumn{2}{@{}l@{}}{\large\bfseries I. Endliche Wörter}\\*[.8ex]

\textbf{Wörter als kleinste abgeschlossene Menge}
\BXXVIIBlickFormel{%
  \begin{aligned}
  \mathcal E_A&\coloneqq\FinSubsets{\mathbb N\times A},\\[-2pt]
  J(w,a)&\coloneqq w\cup\{(\FiniteCard(w),a)\},\\[-2pt]
  \mathfrak C_A&\coloneqq
    \bigl\{C\in\powerset(\mathcal E_A)\bigm|{}\\[-2pt]
    &\quad\varnothing\in C\ \land{}\\[-2pt]
    &\quad\forall w\in C\,\forall a\in A:\\[-2pt]
    &\qquad J(w,a)\in C\bigr\},\\[-2pt]
  \FiniteWords A&\coloneqq\bigcap\mathfrak C_A.
  \end{aligned}}
Dabei sind \(w\in \mathcal E_A\) und \(a\in A\) die Voraussetzungen für
\(J(w,a)\). Der Umgebungsbereich \(\mathcal E_A\) enthält alle endlichen
Teilmengen; erst Anfang und Anfügung bestimmen die Wörter.
\BXXVIIBlickQuellen{%
  \FormulaRefAuto{FiniteOccurrenceSetUniverseDef}[def];
  \FormulaRefAuto{OccurrenceCardinalAdjunctionDef}[def];
  \FormulaRefAuto{WordGenerationFamilyDef}[def];
  \FormulaRefAuto{FiniteWordSetDef}[def]}
&
\textit{Bewiesener Anschluss an endliche Folgen}
\BXXVIIBlickFormel{%
  \begin{aligned}
  w\in\FiniteWords A\ \eqvdash\ {}&
    w\in\powerset(\mathbb N\times A)\ \land{}\\[-2pt]
  &\exists n\in\mathbb N\;(w\colon\NatSeg n\to A),\\[2pt]
  w\in\FiniteWords A\ \vdash\ {}&
    \exists!n\in\mathbb N\;(w\colon\NatSeg n\to A).
  \end{aligned}}
\BXXVIIBlickQuellen{%
  \FormulaRefAuto{FiniteWordSetMembership}[thm];
  \FormulaRefAuto{FiniteWordLengthUnique}[thm]}
Anfang und Anfügungsschritt liefern die Wortinduktion. Die lokale Regel
fasst beide Teilbeweise mit der Wortzugehörigkeit zusammen.
\BXXVIIBlickQuellen{\FormulaRefAuto{GeneratedWordInductionRule}[thm]}\\
\addlinespace[.8ex]

\textbf{Wortlänge}
\BXXVIIBlickFormel{%
  \begin{aligned}
  w\in\FiniteWords A\ \vdash\quad&\\[-7pt]
  \WordLength w\coloneqq{}&\FiniteCard(w).
  \end{aligned}}
Die Länge zählt die aufeinanderfolgenden Positionen und damit die
Anfügungsschritte. Grundlage ist die endliche Kardinalität aus Band~20.
\BXXVIIBlickQuellen{\FormulaRefAuto{FiniteWordLengthDef}[def];
  \FormulaRefAuto{FiniteWordLengthCardinality}[thm]}
&
\textit{Typisierung und Identifikation}
\BXXVIIBlickFormel{%
  \begin{aligned}
  w\in\FiniteWords A\ \vdash\quad
    &\WordLength w\in\mathbb N,\\[-2pt]
    &w\colon\NatSeg{\WordLength w}\to A,\\[2pt]
  n\in\mathbb N\dsep w\colon\NatSeg n\to A\ \vdash\quad
    &\WordLength w=n.
  \end{aligned}}
\BXXVIIBlickQuellen{%
  \FormulaRefAuto{FiniteWordTyping}[thm];
  \FormulaRefAuto{FiniteWordLengthFromTyping}[thm]}\\
\addlinespace[.8ex]

\textbf{Leerwort und Nichtleerwörter}
\BXXVIIBlickFormel{%
  \EmptyWord A\coloneqq\varnothing,
  \qquad
  \NonemptyWords A\coloneqq
    \FiniteWords A\setminus\{\EmptyWord A\}.}
\BXXVIIBlickQuellen{%
  \FormulaRefAuto{EmptyWordDef}[def];
  \FormulaRefAuto{NonemptyWordSetDef}[def]}
&
\textit{Länge und Elementkriterium}
\BXXVIIBlickFormel{%
  \begin{aligned}
  &\EmptyWord A\in\FiniteWords A,
    \qquad \WordLength{\EmptyWord A}=0,\\[2pt]
  &w\in\NonemptyWords A\ \eqvdash\quad
    w\in\FiniteWords A\land\WordLength w\neq0.
  \end{aligned}}
\BXXVIIBlickQuellen{%
  \FormulaRefAuto{EmptyWordLength}[thm];
  \FormulaRefAuto{NonemptyWordMembership}[thm]}\\
\addlinespace[.8ex]

\textbf{Buchstabenwort}
\BXXVIIBlickFormel{%
  a\in A\ \vdash\quad
  \LetterWord a\coloneqq\{(0,a)\}.}
\BXXVIIBlickQuellen{\FormulaRefAuto{LetterWordDef}[def]}
&
\textit{Grunddaten}
\BXXVIIBlickFormel{%
  a\in A\ \vdash\quad
  \begin{aligned}
  &\LetterWord a\colon\NatSeg 1\to A,\\[-2pt]
  &\LetterWord a\in\NonemptyWords A,\\[-2pt]
  &\WordLength{\LetterWord a}=1,
    \qquad \LetterWord a(0)=a.
  \end{aligned}}
\BXXVIIBlickQuellen{\FormulaRefAuto{LetterWordFacts}[thm]}\\
\midrule

\textbf{Primitive Anfügung eines Buchstabens}
\BXXVIIBlickFormel{%
  \sigma_A(u,a)=u\cup\{(\WordLength u,a)\}.}
Die Anfügung wird vor der allgemeinen Konkatenation eingeführt.
\BXXVIIBlickQuellen{\FormulaRefAuto{WordModelSuccessorDef}[def]}
&
\textit{Das konkrete Modell erfüllt W0 bis W3}
Die Minimalität folgt bereits aus der Konstruktion als kleinste
abgeschlossene Menge. Der Anschluss an endliche Folgen liefert
Typisierung, Trennung vom Leerwort und gemeinsame Injektivität.
\BXXVIIBlickQuellen{%
  \FormulaRefAuto{EarlyWordModelMinimality}[thm];
  \FormulaRefAuto{WordStructureConcreteModel}[thm]}\\
\addlinespace[.8ex]

\textbf{Frühe Wortaxiome}
\BXXVIIBlickFormel{%
  e\in W,\qquad s\colon W\times A\to W.}
W1: Der Anfang ist kein Nachfolger. W2: Die Anfügung ist gemeinsam
injektiv. W3: Jede Teilmenge mit Anfang und Anfügungsabschluss ist
der ganze Träger.
\BXXVIIBlickQuellen{\FormulaRefAuto{WordStructureDef}[def]}
&
\textit{Zerlegung und Induktion aus den Axiomen}
Jedes Wort ist der Anfang oder eindeutig ein Nachfolger.
Eine Eigenschaft, die für den Anfang gilt und bei Anfügung erhalten
bleibt, gilt für sämtliche Wörter der Struktur.
\BXXVIIBlickQuellen{%
  \FormulaRefAuto{WordStructureDecomposition}[thm];
  \FormulaRefAuto{WordStructureInduction}[thm];
  \FormulaRefAuto{WordStructureInductionRule}[thm]}
Die lokale Regel gilt in jeder Wortstruktur und damit auch im konkreten Modell.\\
\addlinespace[.8ex]

\textbf{Rekursionsdaten einer Wortstruktur}
\BXXVIIBlickFormel{%
  x_0\in X,\qquad r\colon X\times A\to X.}
Die Rekursionsfunktion wird als kleinster abgeschlossener Graph
unmittelbar aus W0 bis W3 gewonnen.
&
\textit{Axiomatische Rekursion und konkreter Modellfall}
\BXXVIIBlickFormel{%
  \begin{aligned}
  &\exists!F\colon W\to X:\\[-2pt]
  &F(e)=x_0,\\[-2pt]
  &F(s(u,a))=r(F(u),a).
  \end{aligned}}
\BXXVIIBlickQuellen{%
  \FormulaRefAuto{WordStructureRecursion}[thm];
  \FormulaRefAuto{FiniteWordRecursion}[thm]}\\
\midrule

\textbf{Konkatenation}
\BXXVIIBlickFormel{%
  \begin{aligned}
  u,v\in\FiniteWords A\ \vdash\quad&\\[-7pt]
  \WordConcat uv&\coloneqq u\cup\\[-2pt]
  &\quad\{(\WordLength u+j,v(j))\mid{}\\[-2pt]
  &\qquad j\in\NatSeg{\WordLength v}\}.
  \end{aligned}}
\BXXVIIBlickQuellen{\FormulaRefAuto{WordConcatenationDef}[def]}
&
\textit{Typisierung und Koordinaten}
\BXXVIIBlickFormel{%
  \begin{aligned}
  &\WordConcat uv\colon
    \NatSeg{\WordLength u+\WordLength v}\to A,\\[-2pt]
  &i\in\NatSeg{\WordLength u}\ \vdash\quad
    \WordConcat uv(i)=u(i),\\[-2pt]
  &j\in\NatSeg{\WordLength v}\ \vdash\quad
    \WordConcat uv(\WordLength u+j)=v(j).
  \end{aligned}}
\BXXVIIBlickQuellen{\FormulaRefAuto{WordConcatenationTyping}[thm]}\\
\addlinespace[.8ex]

\textit{Die zweite Wortkoordinate wird also um
\(\WordLength u\) verschoben.}
&
\textit{Länge und Rechengesetze}
\BXXVIIBlickFormel{%
  \begin{aligned}
  &\WordLength{\WordConcat uv}=\WordLength u+\WordLength v,\\[-2pt]
  &\WordConcat{\EmptyWord A}w=w=
    \WordConcat w{\EmptyWord A},\\[-2pt]
  &\WordConcat{\WordConcat uv}w=
    \WordConcat u{\WordConcat vw}.
  \end{aligned}}
\BXXVIIBlickQuellen{%
  \FormulaRefAuto{WordConcatenationLength}[thm];
  \FormulaRefAuto{WordConcatenationIdentity}[thm];
  \FormulaRefAuto{WordConcatenationAssociative}[thm]}
Die Rechengesetze werden mit der aus W3 gewonnenen Wortinduktion
bewiesen.\\
\addlinespace[.8ex]

\textit{Fester Schnitt}
\BXXVIIBlickFormel{%
  \WordLength u=\WordLength{u'}\dsep
  \WordConcat uv=\WordConcat{u'}{v'}.}
&
\textit{Kürzung am festgelegten Schnitt}
\BXXVIIBlickFormel{%
  u=u',\qquad v=v'.}
\BXXVIIBlickQuellen{%
  \FormulaRefAuto{WordConcatenationFixedCutCancellation}[thm]}\\
\midrule

\textbf{Konstruktoren nichtleerer Wörter}
\BXXVIIBlickFormel{%
  a\longmapsto\LetterWord a,
  \qquad
  (u,a)\longmapsto\WordConcat u{\LetterWord a}.}
&
\textit{Eindeutige Rechtszerlegung}
\BXXVIIBlickFormel{%
  \begin{aligned}
  w\in\NonemptyWords A\ \vdash\quad
  \exists!u\in\FiniteWords A\;
  \exists!a\in A\;{}\\[-2pt]
  w=\WordConcat u{\LetterWord a}.
  \end{aligned}}
\BXXVIIBlickQuellen{%
  \FormulaRefAuto{NonemptyWordRightDecomposition}[thm]}
\textit{Dabei darf der Präfix \(u\) das Leerwort sein.}\\
\addlinespace[.8ex]

\textbf{Rumpf und Endbuchstabe}
\BXXVIIBlickFormel{%
  \begin{aligned}
  &\WordRumpfMap A\colon\NonemptyWords A\to\FiniteWords A,\\[-2pt]
  &\WordEndMap A\colon\NonemptyWords A\to A.
  \end{aligned}}
\BXXVIIBlickQuellen{%
  \FormulaRefAuto{WordRumpfDef}[def];
  \FormulaRefAuto{WordEndDef}[def]}
&
\textit{Zerlegen und wieder zusammensetzen}
\BXXVIIBlickFormel{%
  \begin{aligned}
  &w\in\NonemptyWords A\ \vdash\\[-2pt]
  &\quad w=\WordConcat{\WordRumpf A w}{\LetterWord{\WordEnd A w}},\\[2pt]
  &u\in\FiniteWords A\dsep a\in A\ \vdash\\[-2pt]
  &\quad\WordRumpf A{\WordConcat u{\LetterWord a}}=u,\\[-2pt]
  &\quad\WordEnd A{\WordConcat u{\LetterWord a}}=a.
  \end{aligned}}
\BXXVIIBlickQuellen{%
  \FormulaRefAuto{WordRightDecompositionEquation}[thm];
  \FormulaRefAuto{WordRumpfAppend}[thm];
  \FormulaRefAuto{WordEndAppend}[thm]}\\
\addlinespace[.8ex]

\textbf{Induktionsdaten}
\BXXVIIBlickFormel{%
  \begin{aligned}
  &\forall a\in A\;P(\LetterWord a),\\[-2pt]
  &\forall u\in\NonemptyWords A\;\forall a\in A:\\[-2pt]
  &\qquad P(u)\to P(\WordConcat u{\LetterWord a}).
  \end{aligned}}
&
\textit{Induktion über nichtleere Wörter}
\BXXVIIBlickFormel{%
  \forall w\in\NonemptyWords A\;P(w).}
\BXXVIIBlickQuellen{\FormulaRefAuto{NonemptyWordInduction}[thm]}\\
\addlinespace[.8ex]

\textbf{Rekursionsdaten}
\BXXVIIBlickFormel{%
  f\colon A\to X,
  \qquad s\colon X\times A\to X.}
&
\textit{Rekursion über nichtleere Wörter}
\BXXVIIBlickFormel{%
  \begin{aligned}
  &\exists!F\colon\NonemptyWords A\to X:\\[-2pt]
  &F(\LetterWord a)=f(a),\\[-2pt]
  &F(\WordConcat u{\LetterWord a})=s(F(u),a).
  \end{aligned}}
\BXXVIIBlickQuellen{\FormulaRefAuto{NonemptyWordRecursion}[thm]}
Sie folgt aus der zuvor bewiesenen Rekursion für alle Wörter.\\
\addlinespace[.8ex]

\textbf{Linksfaltung}
\BXXVIIBlickFormel{%
  \begin{aligned}
  &\mathsf{BinOp}(A,\star)\ \vdash\\[-2pt]
  &\Pi_\star\coloneqq
    \iota F\colon\NonemptyWords A\to A,\\[-2pt]
  &\forall a\in A\quad F(\LetterWord a)=a,\\[-2pt]
  &\forall u\in\NonemptyWords A\;\forall a\in A:\\[-2pt]
  &\qquad F(\WordConcat u{\LetterWord a})=F(u)\star a.
  \end{aligned}}
\BXXVIIBlickQuellen{\FormulaRefAuto{LeftWordFoldDef}[def]}
&
\textit{Rekursionsgleichungen}
\BXXVIIBlickFormel{%
  \begin{aligned}
  &\WordFold\star{\LetterWord a}=a,\\[-2pt]
  &\WordFold\star{\WordConcat u{\LetterWord a}}
    =\WordFold\star u\star a.
  \end{aligned}}
\BXXVIIBlickQuellen{%
  \FormulaRefAuto{LeftWordFoldFunction}[thm];
  \FormulaRefAuto{LeftWordFoldEquations}[thm]}
\textit{Ohne ausgezeichnetes neutrales Element wird das Leerwort nicht
gefaltet.}\\
\midrule

\multicolumn{2}{@{}l@{}}{%
  \large\bfseries II. Wortcodes und ihre vollen planaren Positionsbäume}\\*[.8ex]

\textbf{Baumalphabet und Blattcode}
\BXXVIIBlickFormel{%
  \begin{aligned}
  \TreeAlphabet A&\coloneqq
    (\{0\}\times A)\cup\\[-2pt]
    &\qquad(\{1\}\times\mathbb N),\\[-2pt]
  a\in A\ \vdash\quad
  \LeafTree A a&\coloneqq\LetterWord{(0,a)}.
  \end{aligned}}
\BXXVIIBlickQuellen{%
  \FormulaRefAuto{TreeTokenAlphabetDef}[def];
  \FormulaRefAuto{TreeLeafCodeDef}[def]}
&
\textit{Blattinjektivität}
\BXXVIIBlickFormel{%
  a,b\in A\ \vdash\quad
  \bigl(\LeafTree A a=\LeafTree A b\ \leftrightarrow\ a=b\bigr).}
\BXXVIIBlickQuellen{\FormulaRefAuto{TreeConstructorNoConfusion}[thm]}\\
\addlinespace[.8ex]

\textbf{Knotencode}
\BXXVIIBlickFormel{%
  \begin{aligned}
  &S,T\in\NonemptyWords{\TreeAlphabet A}\ \vdash\\[-2pt]
  &\NodeTree A S T\coloneqq
    \LetterWord{(1,\WordLength S)}\frown\\[-2pt]
  &\hspace{7em}(\WordConcat ST).
  \end{aligned}}
\BXXVIIBlickQuellen{\FormulaRefAuto{TreeNodeCodeDef}[def]}
&
\textit{Konstruktortrennung und Knoteninjektivität}
\BXXVIIBlickFormel{%
  \begin{aligned}
  &\LeafTree A a\neq\NodeTree A S T,\\[-2pt]
  &\NodeTree A S T=\NodeTree A{S'}{T'}\ \vdash\\[-2pt]
  &\qquad S=S',\quad T=T'.
  \end{aligned}}
\BXXVIIBlickQuellen{\FormulaRefAuto{TreeConstructorNoConfusion}[thm]}
\textit{Die gespeicherte Länge \(\WordLength S\) legt den Schnitt des
Knotencodes fest.}\\
\midrule

\textbf{Erzeugungsschritt}
\BXXVIIBlickFormel{%
  \begin{aligned}
  &\TreeClosure A X\coloneqq\\[-2pt]
  &\quad\{\LeafTree A a\mid a\in A\}\cup\\[-2pt]
  &\quad\{\NodeTree A S T\mid S,T\in X\}.
  \end{aligned}}
\BXXVIIBlickQuellen{\FormulaRefAuto{TreeGenerationStepDef}[def]}
&
\textit{Abbildung und Monotonie}
\BXXVIIBlickFormel{%
  \begin{aligned}
  &\mathcal C_A\colon
    \powerset(\NonemptyWords{\TreeAlphabet A})
    \to\powerset(\NonemptyWords{\TreeAlphabet A}),\\[-2pt]
  &X\subseteq Y\ \vdash\quad
    \TreeClosure A X\subseteq\TreeClosure A Y.
  \end{aligned}}
\BXXVIIBlickQuellen{%
  \FormulaRefAuto{TreeGenerationStepFunction}[thm];
  \FormulaRefAuto{TreeGenerationStepMonotonicity}[thm]}\\
\addlinespace[.8ex]

\textbf{Baumstufen}
\BXXVIIBlickFormel{%
  n\in\mathbb N\ \vdash\quad
  \TreeStage A n\coloneqq
  \RecFun{\varnothing}{\mathcal C_A}(n).}
\BXXVIIBlickQuellen{\FormulaRefAuto{TreeStageDef}[def]}
&
\textit{Stufengleichungen und Verschachtelung}
\BXXVIIBlickFormel{%
  \begin{aligned}
  &\TreeStage A0=\varnothing,\\[-2pt]
  &n\in\mathbb N\ \vdash\quad
    \TreeStage A{n+1}=\TreeClosure A{\TreeStage A n},\\[-2pt]
  &m,n\in\mathbb N\dsep m\leq n\ \vdash\quad
    \TreeStage A m\subseteq\TreeStage A n.
  \end{aligned}}
\BXXVIIBlickQuellen{%
  \FormulaRefAuto{TreeStageZeroEquation}[thm];
  \FormulaRefAuto{TreeStageSuccessorEquation}[thm];
  \FormulaRefAuto{TreeStageNested}[thm]}\\
\addlinespace[.8ex]

\textbf{Menge der Klammerungsbäume}
\BXXVIIBlickFormel{%
  \begin{aligned}
  \BracketTreeSet A\coloneqq
  \{R\in\NonemptyWords{\TreeAlphabet A}\mid{}&
    \exists n\in\mathbb N\;{}\\[-2pt]
    &R\in\TreeStage A n\}.
  \end{aligned}}
\BXXVIIBlickQuellen{\FormulaRefAuto{FullPlanarBinaryTreeSetDef}[def]}
&
\textit{Das konkrete Modell erfüllt B0 bis B4}
\BXXVIIBlickFormel{%
  \begin{aligned}
  &\lambda_A\colon A\to\BracketTreeSet A,\\[-2pt]
  &\kappa_A\colon\BracketTreeSet A\times\BracketTreeSet A
      \to\BracketTreeSet A.
  \end{aligned}}
Die Konstruktoren sind injektiv und haben getrennte Bilder.
Die Minimalität folgt durch natürliche Induktion über die Baumstufen.
\BXXVIIBlickQuellen{%
  \FormulaRefAuto{TreeCodeMinimality}[thm];
  \FormulaRefAuto{BinaryTreeStructureConcreteModel}[thm]}\\
\addlinespace[.8ex]

\textbf{Frühe Baumaxiome}
\BXXVIIBlickFormel{%
  \ell\colon A\to T,\qquad k\colon T\times T\to T.}
B1 und B2: Beide Konstruktoren sind injektiv.
B3: Ihre Bilder sind getrennt.
B4: Jede Teilmenge mit allen Blättern und Knotenabschluss ist
der ganze Träger.
\BXXVIIBlickQuellen{\FormulaRefAuto{BinaryTreeStructureDef}[def]}
&
\textit{Zerlegung und Induktion aus den Axiomen}
Jeder Baum ist eindeutig ein beschriftetes Blatt oder ein Knoten
mit zwei geordneten Teilbäumen. Die Minimalität liefert die
Strukturinduktion für jede solche Baumstruktur.
\BXXVIIBlickQuellen{%
  \FormulaRefAuto{BinaryTreeStructureInduction}[thm];
  \FormulaRefAuto{TreeConstructorDecomposition}[thm]}\\
\midrule

\textbf{Daten der Strukturinduktion}
\BXXVIIBlickFormel{%
  \begin{aligned}
  &\forall a\in A\;P(\LeafTree A a),\\[-2pt]
  &\forall S,T\in\BracketTreeSet A:\\[-2pt]
  &\qquad P(S)\land P(T)\to\\[-2pt]
  &\qquad P(\NodeTree A S T).
  \end{aligned}}
&
\textit{Strukturinduktion}
\BXXVIIBlickFormel{%
  \forall R\in\BracketTreeSet A\;P(R).}
\BXXVIIBlickQuellen{%
  \FormulaRefAuto{FullPlanarBinaryTreeStructuralInduction}[thm]}\\
\addlinespace[.8ex]

\textbf{Daten der Strukturrekursion}
\BXXVIIBlickFormel{%
  f\colon A\to X,
  \qquad g\colon X\times X\to X.}
&
\textit{Axiomatische Rekursion und konkreter Modellfall}
\BXXVIIBlickFormel{%
  \begin{aligned}
  &\exists!E\colon\BracketTreeSet A\to X:\\[-2pt]
  &E(\LeafTree A a)=f(a),\\[-2pt]
  &E(\NodeTree A S T)=g(E(S),E(T)).
  \end{aligned}}
\BXXVIIBlickQuellen{%
  \FormulaRefAuto{BinaryTreeStructureRecursion}[thm];
  \FormulaRefAuto{FullPlanarBinaryTreeRecursion}[thm]}
Der allgemeine Satz folgt direkt aus B0 bis B4 durch einen kleinsten
abgeschlossenen Graphen; der angezeigte Satz ist seine Anwendung
auf das konkrete Modell.\\
\addlinespace[.8ex]

\textbf{Der benannte Baumrekursor}
\BXXVIIBlickFormel{%
  \begin{aligned}
  &f\colon A\to X,\quad g\colon X\times X\to X,\\[-2pt]
  &R\coloneqq\TreeRecursor A X f g.
  \end{aligned}}
\textit{Blätter durch Werte ersetzen, Teilbaumwerte verknüpfen.}
\BXXVIIBlickQuellen{\FormulaRefAuto{TreeRecursorDef}[def]}
&
\textit{Einheitliche Rechen- und Vergleichsregeln}
\BXXVIIBlickFormel{%
  \begin{aligned}
  &R\colon\BracketTreeSet A\to X,\\[-2pt]
  &a\in A\ \vdash\ R(\LeafTree A a)=f(a),\\[-2pt]
  &S,T\in\BracketTreeSet A\ \vdash\\[-2pt]
  &\quad R(\NodeTree A S T)=g(R(S),R(T)).
  \end{aligned}}
\textit{Eine gleich typisierte Funktion mit denselben Blatt- und
Knotenregeln ist gleich \(R\).}
\BXXVIIBlickQuellen{%
  \FormulaRefAuto{TreeRecursorEquations}[thm];
  \FormulaRefAuto{TreeRecursorIdentification}[thm]}\\
\addlinespace[.8ex]

\textbf{Positionsgraph eines Baumcodes}
\BXXVIIBlickFormel{%
  \begin{aligned}
  \BinaryAddresses&\coloneqq\FiniteWords{\{0,1\}},\\[-2pt]
  P_R&\coloneqq\TreePositions A R,\\[-2pt]
  \TreePositionGraph A R&\coloneqq
    \bigl(P_R,\AddressEdges{P_R},
      \EmptyWord{\{0,1\}},\\[-2pt]
  &\hspace{6em}\AddressLeft{P_R},\AddressRight{P_R}\bigr).
  \end{aligned}}
\BXXVIIBlickQuellen{%
  \FormulaRefAuto{TreePositionGraphDef}[def]}
&
\textit{Der Wortcode trägt einen besonderen Baum}
\BXXVIIBlickFormel{%
  \begin{aligned}
  P_R&\coloneqq\TreePositions A R,\\[-2pt]
  R\in\BracketTreeSet A\ \vdash\quad
    &\Finite{P_R},\\[-2pt]
    &\TreeStruct{P_R,\AddressEdges{P_R}}.
  \end{aligned}}
\BXXVIIBlickQuellen{%
  \FormulaRefAuto{BracketCodePositionGraphIsFiniteOrderedFullBinaryTree}[thm]}
\textit{Der Positionsgraph ist sogar ein geordneter voller Binärbaum.
Seine Knoten sind Adressen; daher bleiben auch gleiche Teilbaumcodes an
verschiedenen Positionen verschieden.}\\
\midrule

\multicolumn{2}{@{}l@{}}{%
  \large\bfseries III. Blattwörter, Klammerungen und Auswertung}\\*[.8ex]

\textbf{Einbuchstabenabbildung und Blattwort}
\BXXVIIBlickFormel{%
  \begin{aligned}
  &\operatorname{ein}_A\colon A\to\NonemptyWords A,
    \qquad \operatorname{ein}_A(a)\coloneqq\LetterWord a,\\[2pt]
  &\operatorname{wort}_A\coloneqq\\[-2pt]
  &\quad\TreeRecursor A{\NonemptyWords A}{\operatorname{ein}_A}
    {\BinaryOpFunction{\NonemptyWords A}{\frown}}.
  \end{aligned}}
\BXXVIIBlickQuellen{%
  \FormulaRefAuto{LetterWordFunctionDef}[def];
  \FormulaRefAuto{TreeLeafWordDef}[def]}
&
\textit{Rekursionsgleichungen des Blattwortes}
\BXXVIIBlickFormel{%
  \begin{aligned}
  &\operatorname{wort}_A\colon
    \BracketTreeSet A\to\NonemptyWords A,\\[-2pt]
  &\TreeLeafWord A{\LeafTree A a}=\LetterWord a,\\[-2pt]
  &\TreeLeafWord A{\NodeTree A S T}=\\[-2pt]
  &\qquad\WordConcat{\TreeLeafWord A S}{\TreeLeafWord A T}.
  \end{aligned}}
\BXXVIIBlickQuellen{\FormulaRefAuto{TreeLeafWordEquations}[thm]}\\
\addlinespace[.8ex]

\textbf{Ein Baum klammert ein Wort}
\BXXVIIBlickFormel{%
  \TreeBrackets A T w\coloneqq
    (T,w)\in\operatorname{wort}_A.}
\BXXVIIBlickQuellen{\FormulaRefAuto{TreeBracketingRelationDef}[def]}
\textit{Die Relation enthält zugleich die Baum- und Worttypisierung.}
&
\textit{Blattregel und Knotenregel}
\BXXVIIBlickFormel{%
  \begin{aligned}
  &a\in A\ \vdash\\[-2pt]
  &\quad\TreeBrackets A{\LeafTree A a}{\LetterWord a},\\[2pt]
  &\TreeBrackets A S u\dsep\TreeBrackets A T v\ \vdash\\[-2pt]
  &\quad\TreeBrackets A{\NodeTree A S T}{\WordConcat u v}.
  \end{aligned}}
\BXXVIIBlickQuellen{%
  \FormulaRefAuto{TreeBracketingLeafIntroduction}[thm];
  \FormulaRefAuto{TreeBracketingNodeIntroduction}[thm]}\\
\addlinespace[.8ex]

\textbf{Ein Blatt rechts anfügen}
\BXXVIIBlickFormel{%
  \begin{aligned}
  &T\in\BracketTreeSet A\dsep a\in A\ \vdash\\[-2pt]
  &\TreeAppendLeaf A T a\coloneqq\\[-2pt]
  &\qquad\NodeTree A T{\LeafTree A a}.
  \end{aligned}}
\BXXVIIBlickQuellen{\FormulaRefAuto{TreeAppendLeafDef}[def]}
&
\textit{Baum und Wort gemeinsam erweitern}
\BXXVIIBlickFormel{%
  \begin{aligned}
  &\TreeBrackets A T u\dsep a\in A\ \vdash\\[-2pt]
  &\quad\TreeBrackets A{\TreeAppendLeaf A T a}
    {\WordConcat u{\LetterWord a}},\\[2pt]
  &w\in\NonemptyWords A\ \vdash
    \exists T\,(\TreeBrackets A T w).
  \end{aligned}}
\BXXVIIBlickQuellen{%
  \FormulaRefAuto{TreeBracketingAppendLeafIntroduction}[thm];
  \FormulaRefAuto{TreeBracketingExistence}[thm]}\\
\addlinespace[.8ex]

\textbf{Der Linksbaum eines Wortes}
\BXXVIIBlickFormel{%
  \begin{aligned}
  &\LeftBracketingMap A\colon\NonemptyWords A\to\BracketTreeSet A,\\[-2pt]
  &a\in A\ \vdash\\[-2pt]
  &\quad\LeftBracketing A{\LetterWord a}=\LeafTree A a,\\[-2pt]
  &u\in\NonemptyWords A\dsep a\in A\ \vdash\\[-2pt]
  &\quad\LeftBracketing A{\WordConcat u{\LetterWord a}}\\[-2pt]
  &\qquad=\TreeAppendLeaf A{\LeftBracketing A u}a.
  \end{aligned}}
\BXXVIIBlickQuellen{%
  \FormulaRefAuto{LeftBracketingDef}[def];
  \FormulaRefAuto{LeftBracketingLetterEquation}[thm];
  \FormulaRefAuto{LeftBracketingAppendEquation}[thm]}
&
\textit{Eine konkrete Klammerung und ihre Auswertung}
\BXXVIIBlickFormel{%
  \begin{aligned}
  &w\in\NonemptyWords A\ \vdash\\[-2pt]
  &\quad\TreeBrackets A{\LeftBracketing A w}w,\\[2pt]
  &\mathsf{BinOp}(A,\star)\dsep w\in\NonemptyWords A\ \vdash\\[-2pt]
  &\quad\TreeEvaluation\star{\LeftBracketing A w}=\WordFold\star w.
  \end{aligned}}
\textit{Den Linksbaum auswerten ergibt die Linksfaltung.}
\BXXVIIBlickQuellen{%
  \FormulaRefAuto{LeftBracketingProperty}[thm];
  \FormulaRefAuto{LeftBracketingEvaluation}[thm]}\\
\addlinespace[.8ex]

\textbf{Klammerungsmenge eines Nichtleerwortes}
\BXXVIIBlickFormel{%
  \begin{aligned}
  &w\in\NonemptyWords A\ \vdash\\[-2pt]
  &\WordBracketings A w\coloneqq\\[-2pt]
  &\quad\{T\in\BracketTreeSet A\mid
    \TreeLeafWord A T=w\}.
  \end{aligned}}
\BXXVIIBlickQuellen{\FormulaRefAuto{WordBracketingSetDef}[def]}
&
\textit{Elementkriterium und Existenz}
\BXXVIIBlickFormel{%
  \begin{aligned}
  &w\in\NonemptyWords A\ \vdash\\[-2pt]
  &\bigl(T\in\WordBracketings A w\ \leftrightarrow\\[-2pt]
  &\quad(T\in\BracketTreeSet A\ \land
    \TreeLeafWord A T=w)\bigr),\\[2pt]
  &w\in\NonemptyWords A\ \vdash
    \WordBracketings A w\neq\varnothing.
  \end{aligned}}
\BXXVIIBlickQuellen{%
  \FormulaRefAuto{WordBracketingMembership}[thm];
  \FormulaRefAuto{WordBracketingExistence}[thm]}\\
\addlinespace[.8ex]

\textbf{Auswertung eines Klammerungsbaums}
\BXXVIIBlickFormel{%
  \begin{aligned}
  &\mathsf{BinOp}(A,\star)\ \vdash\\[-2pt]
  &\operatorname{ev}_\star\coloneqq\\[-2pt]
  &\quad\TreeRecursor A A{\Id_A}{\BinaryOpFunction A\star}.
  \end{aligned}}
\BXXVIIBlickQuellen{\FormulaRefAuto{TreeEvaluationDef}[def]}
&
\textit{Typisierung und Rekursionsgleichungen}
\BXXVIIBlickFormel{%
  \begin{aligned}
  &\operatorname{ev}_\star\colon\BracketTreeSet A\to A,\\[-2pt]
  &\TreeEvaluation\star{\LeafTree A a}=a,\\[-2pt]
  &\TreeEvaluation\star{\NodeTree A S T}=
    \TreeEvaluation\star S\star\TreeEvaluation\star T.
  \end{aligned}}
\BXXVIIBlickQuellen{\FormulaRefAuto{TreeEvaluationEquations}[thm]}
\textit{Erst die Assoziativität in Band~28 macht die Auswertung von der
gewählten Klammerung unabhängig.}\\

\end{longtable}
\endgroup

\section*{Vom konkreten Modell zu frühen Strukturaxiomen}

Das Anfügungsmodell beginnt im Bereich aller endlichen Teilmengen von
\(\mathbb N\times A\). Die Wortmenge wird als Schnitt aller Teilmengen
dieses Bereichs definiert, die \(\varnothing\) enthalten und unter
\(u\mapsto u\cup\{(\FiniteCard(u),a)\}\) für \(a\in A\) abgeschlossen sind
(\FormulaRefAuto{FiniteWordSetDef}[def]). Wie beim Aufbau der natürlichen
Zahlen durch eine kleinste induktive Menge folgt die Minimalität
unmittelbar aus dieser Konstruktion. Die endlichen Teilmengen und ihre
Kardinalzahlen liefert Band~20; eine Wortmenge wird dabei noch nicht
vorausgesetzt.

Anschließend wird bewiesen, dass genau die nullbasierten endlichen Folgen
aus Band~21 entstehen. Dieser Anschluss begründet die Positionsdarstellung
und die eindeutige Zerlegung. Auf der Wortmenge wird der Erzeugungsschritt
als Anfügungsfunktion \(\sigma_A\) gefasst, noch vor der Konkatenation.
\FormulaRefAuto{WordStructureDef}[def] bündelt die so bewiesenen
Eigenschaften als W0 bis W3 für beliebige Wortstrukturen.

Beim Baumaufbau werden zunächst die Konstruktoren und die endlichen
Erzeugungsstufen untersucht. Natürliche Induktion über diese Stufen
beweist die Minimalität des konkreten Baumträgers
(\FormulaRefAuto{TreeCodeMinimality}[thm]).
Danach führt \FormulaRefAuto{BinaryTreeStructureDef}[def]
die Bedingungen B0 bis B4 ein: typisierte, injektive Blatt- und
Knotenkonstruktoren mit getrennten Bildern sowie Minimalität.
Die beiden Axiomsysteme sind Bedingungen an Strukturen; ihre bereits
konstruierten Modelle erfordern keine zusätzlichen Mengenexistenzaxiome.

\widowpenalty=10000
Aus den frühen Axiomen folgen eindeutige Zerlegung und Induktion.
Die Rekursion wird unmittelbar durch den Schnitt aller unter den
vorgegebenen Rekursionsregeln abgeschlossenen Relationsgraphen bewiesen;
siehe \FormulaRefAuto{WordStructureRecursion}[thm] und
\FormulaRefAuto{BinaryTreeStructureRecursion}[thm].
Erst daraus folgen die konkreten Rekursionssätze und die Eindeutigkeit
der Strukturen bis auf einen eindeutig bestimmten strukturtreuen
Isomorphismus. Die späteren Wortzerlegungen, Induktionsbeweise,
Konkatenationsgesetze und Auswertungen greifen auf diese Ergebnisse
zurück. Die konkrete Darstellung kann deshalb gewechselt werden,
ohne die Aufbauregeln zu verändern.

Der Anschluss an allgemeine Bäume erfolgt über
\FormulaRefAuto{B27RankedParentSystemTree}[thm]: Eine endliche
Knotenmenge mit eindeutigen Eltern und geeigneten Höhenmarken trägt
einen verwurzelten Baum. Die volle geordnete binäre Verzweigung ist
eine zusätzliche Bedingung des Klammerungsmodells.
Für die Syntax in Band~48 dienen Konstruktortrennung, Induktion und
Rekursion als wiederverwendbares Beweismuster. Die dortigen
Konstruktoren unterschiedlicher Stelligkeit benötigen jeweils ihre
eigenen Aufbauregeln; eine binäre Codierung ersetzt diese Regeln nicht.
\endgroup

